\documentclass[12pt,aspectratio=1610,professionalfont,serif,t]{beamer}
\usepackage[osf,p]{newpxtext}
\usepackage{newpxmath}
\usepackage{fontawesome5}

\usepackage{fancyvrb}
\fvset{%frame=single,
commandchars=\\\{\},
%framesep=1mm,
fontfamily=helvetica,
fontsize=\normalsize
}
%\usepackage[bitstream-charter]{mathdesign}
\usepackage{listings}
\lstset{commentstyle=\color{orange}, keywordstyle=\color{yellow} }


\usetheme{Boadilla}
\useoutertheme{split}
\usecolortheme{albatross}


\setbeamertemplate{blocks}[rounded][shadow=false]
\setbeamersize{text margin left=\dimexpr 1cm+1ex\relax, text margin right=\dimexpr 1cm+1ex\relax}
\addtobeamertemplate{headline}{}{\usebeamerfont{frametitle}\vspace*{1ex}}
\setbeamertemplate{navigation symbols}{}
\setbeamertemplate{page number in head/foot}{}
%\beamerdefaultoverlayspecification{<+->}


\setbeamercolor*{structure}{fg=green!60!black,bg=black}
\setbeamercolor*{normal text}{fg=green!60!black,bg=black}
\setbeamercolor{palette primary}{use={structure,normal text},fg=green,bg=structure.bg}
\setbeamercolor{palette secondary}{use={structure,normal text},fg=structure.fg,bg=structure.bg}
\setbeamercolor{palette tertiary}{use={structure,normal text},fg=structure.fg,bg=structure.bg}
\setbeamercolor{palette quaternary}{use={structure,normal text},fg=white,bg=green!30!black}
\setbeamercolor{palette sidebar quaternary}{use={structure,normal text},fg=white,bg=green!30!black}
\setbeamercolor{title in head/foot}{use={structure,normal text},fg=white,bg=green!30!black}

\setbeamercolor*{example text}{use={structure},fg=green!60!black}
\setbeamercolor*{block body}{use={structure},bg=structure.bg}
\setbeamercolor*{block body alerted}{use={structure},bg=structure.bg}
\setbeamercolor*{block body example}{use={structure},bg=structure.bg}
\setbeamercolor*{block title}{parent=structure,bg=structure.bg}
\setbeamercolor*{block title alerted}{use={structure,alerted text},fg=alerted text.fg,bg=structure.bg}
\setbeamertemplate{navigation symbols}{}
\setbeamertemplate{items}[square]
\setbeamercolor{item projected}{fg=white}
\setbeamercolor*{normal text}{use={structure},fg=white!90!blue,bg=structure.bg}
\setbeamercolor*{separation line}{}
\setbeamercolor*{fine separation line}{}
\setbeamercolor{alerted text}{use={structure},fg=structure.fg}

\usepackage[italian]{babel}
\usepackage[utf8]{inputenc}
\usepackage{pgf}
\usepackage{verbatim}
\usepackage{inconsolata}
\usepackage{listings}
\lstset{language=C, frame=single,
	basicstyle=\ttfamily,
	numbers=left, numberstyle=\tiny\color{gray},
	numbersep=5pt, fancyvrb=true
}

\usepackage[noend]{algorithm2e}
\usepackage{pgf}
\usepackage{tikz}
\usepackage{graphicx}
\usetikzlibrary{%
	arrows,
	arrows.meta,
	positioning,
	calc,
	backgrounds,
	chains,
	matrix,
	patterns,
	automata,
	fit,
	graphs,
	decorations,
	decorations.pathmorphing,
	decorations.pathreplacing,
	decorations.markings,
}

\usepgflibrary{shapes,shapes.geometric}


\usepackage{url}


\usepackage{xmpmulti}
\usepackage{transparent}
\graphicspath{{img/}}

\newcommand{\zero}{\ensuremath{\mathtt{0}}\xspace}
\newcommand{\uno}{\ensuremath{\mathtt{1}}\xspace}
\newcommand{\due}{\ensuremath{\mathtt{2}}\xspace}
\newcommand{\tre}{\ensuremath{\mathtt{3}}\xspace}
\newcommand{\kleene}{\ensuremath{{}^{\ast}}\xspace}
\xspaceaddexceptions{\zero \uno \due \tre \kleene}


\newcommand{\simplebox}[1]{
      \begin{tikzpicture}
              \node[text centered, rectangle, draw=black, thick, text width=\linewidth,align=center, rounded corners, minimum height=1em] (C) {#1};
      \end{tikzpicture}
}
%
% In questo file si mettono i comandi utili solo per questo insegnamento.
%
\author{Gianluca Della Vedova}
\title{Linguaggi e Computabilità}
\institute{Univ. Milano--Bicocca\\
	\texttt{https://www.unimib.it/gianluca-della-vedova}}
\date{\today}



\begin{document}


\begin{frame}[fragile]
	\frametitle{Linguaggi e Computabilità: Lezione 12}
\begin{enumerate}
\item
Equivalenza di stati
\end{enumerate}
\end{frame}


\begin{frame}
        \frametitle{Esercizio 1: calcolare quali stati sono equivalenti}
	\begin{center}
		\begin{tikzpicture}[enfa]
			\node[state, initial]   (q0)               {$q_0$};
                        \node[state]            (q1) [below=of q0] {$q_1$};
                        \node[state]            (q1b) [below=of q1] {$q_4$};
			\node[state]            (q2) [right=of q1] {$q_2$};
                        \node[state]            (q2b) [left=of q1b] {$q_6$};
                        \node[state, accepting] (q3b) [right=of q1b] {$q_5$};
                        \node[state, accepting] (q3) [right=of q3b] {$q_3$};
                        % Posizione delle etichette: lo stile enfa usa `auto`, quindi
                        % l'etichetta sta a sinistra rispetto al verso dell'arco.
                        % - node[swap]        sposta l'etichetta dall'altro lato
                        % - node[above], node[below], node[left], node[right]
                        %                     posizione assoluta rispetto all'arco
                        % - node[pos=0.3]     posizione lungo l'arco (0 = sorgente, 1 = destinazione)
                        % - node[near start], node[near end] scorciatoie per pos
                        \path (q0) edge                 node[swap] {\uno}          (q1)
                                   edge [bend left=20]  node       {\zero}          (q2)
                              (q1) edge [bend left=20]    node       {\uno}          (q1b)
                                   edge                 node       {\zero} (q2)
                              (q2) edge [bend left=20]  node       {\uno}         (q3)
                                   edge   node       {\zero}          (q0)
                              (q3) edge   node       {\uno}         (q2)
                                   edge [bend left=-35] node[swap]       {\zero}          (q0)
                              (q1b) edge                 node       {\uno}         (q1)
                                   edge [bend left=-20] node[swap]       {\zero}          (q2b)
                              (q3b) edge                 node       {\uno}         (q2)
                                   edge [bend left=10] node       {\zero}          (q0)
                              (q2b) edge  [bend left=-15]                node[swap]       {\uno}         (q3b)
                                   edge [bend left=10] node       {\zero}          (q0);
		\end{tikzpicture}
	\end{center}
\end{frame}


% \begin{frame}
%         \frametitle{Esercizio 1: automa minimo}
% 	% Soluzione.
% 	% 1. Tutti gli stati sono raggiungibili da q0
% 	%    (q0 -> q1, q2; q1 -> q4; q2 -> q3; q4 -> q6; q6 -> q5).
% 	% 2. Partizione iniziale: {q3, q5} (finali), {q0, q1, q2, q4, q6} (non finali).
% 	% 3. Con 1, q2 -> q3 e q6 -> q5 vanno in uno stato finale,
% 	%    mentre q0 -> q1, q1 -> q4, q4 -> q1 no.
% 	%    Partizione: {q0, q1, q4}, {q2, q6}, {q3, q5}.
% 	% 4. Verifica di stabilità:
% 	%    q0, q1, q4: con 0 in {q2, q6}, con 1 in {q0, q1, q4};
% 	%    q2, q6:     con 0 in {q0, q1, q4}, con 1 in {q3, q5};
% 	%    q3, q5:     con 0 in {q0, q1, q4}, con 1 in {q2, q6}.
% 	%    La partizione è stabile: l'automa minimo ha 3 stati.
% 	\begin{center}
% 		\begin{tikzpicture}[enfa, node distance=2.8cm]
% 			\node[state, initial]   (A)               {$\{q_0, q_1, q_4\}$};
% 			\node[state]            (B) [right=of A] {$\{q_2, q_6\}$};
% 			\node[state, accepting] (C) [right=of B] {$\{q_3, q_5\}$};
% 			\path (A) edge [loop above]   node       {\uno}  ()
% 			          edge [bend left=20] node       {\zero} (B)
% 			      (B) edge [bend left=20] node       {\zero} (A)
% 			          edge [bend left=20] node       {\uno}  (C)
% 			      (C) edge [bend left=20] node       {\uno}  (B)
% 			          edge [bend left=45] node       {\zero} (A);
% 		\end{tikzpicture}
% 	\end{center}
% \end{frame}



\end{document}
